\begin{tikzpicture}[
  font=\scriptsize,
  execute at begin node={\hyphenpenalty=10000},
  tahap/.style={draw, rounded corners, align=center, text width=2.35cm,
    minimum height=1.05cm, fill=blue!8},
  panah/.style={-{Stealth[length=2mm]}, thick}]
  \node[tahap] (n1) at (0,0)
    {Pemicu Pasal 34\\{\tiny\mdseries tujuh keadaan ayat (2)}};
  \node[tahap] (n2) at (3,0)
    {Tujuan dan kebutuhan\\{\tiny\mdseries data apa yang perlu}};
  \node[tahap] (n3) at (6,0)
    {Peta data\\{\tiny\mdseries sumber, penerima, lokasi}};
  \node[tahap] (n4) at (9,0)
    {Dampak pada pasien\\{\tiny\mdseries bukan hanya rupiah}};
  \node[tahap] (n5) at (9,-1.75)
    {Kontrol dan risiko sisa\\{\tiny\mdseries bukti uji sebelum aktif}};
  \node[tahap, fill=orange!12] (n6) at (6,-1.75)
    {Keputusan tujuan\\{\tiny\mdseries direktur pelayanan}};
  \node[tahap, fill=green!10] (n7) at (3,-1.75)
    {Aktif dengan syarat\\{\tiny\mdseries pemilik layanan}};
  \node[tahap] (n8) at (0,-1.75)
    {Pantau dan tinjau ulang\\{\tiny\mdseries bila tujuan atau skala berubah}};
  \node[tahap, fill=red!22] (n9) at (6,-3.5)
    {Hentikan atau naik\\{\tiny\mdseries risiko sisa di atas batas ke direksi}};
  \draw[panah] (n1) -- (n2);
  \draw[panah] (n2) -- (n3);
  \draw[panah] (n3) -- (n4);
  \draw[panah] (n4) -- (n5);
  \draw[panah] (n5) -- (n6);
  \draw[panah] (n6) -- (n7);
  \draw[panah] (n7) -- (n8);
  \draw[panah] (n6) -- node[right, align=left, font=\tiny\mdseries]
    {kurangi data,\\nilai ulang} (n3);
  \draw[panah] (n6) -- (n9);
\end{tikzpicture}
